\subsection{Limit na počet tokenů}
\label{subsec:llm-token-limit}

Tokenizace a kontextové okno byly vysvětleny v \secref[sekci]{subsec:llm-tokens-parameters-context}.
Každý model je omezen maximálním počtem tokenů, které může zpracovat v jednom průchodu inference.
Pokud vstup tento limit překročí, část informací se do modelu vůbec nedostane.
Typicky jsou odříznuty tokeny z konce vstupu, nebo musí být vstup manuálně zkrácen před odesláním.

Pro analýzu zdrojového kódu je toto omezení zvlášť problematické ze dvou důvodů:
\begin{enumerate}
    \item Programátorské chyby mají svou strukturu rozloženou napříč kódem.
    Únik paměti může vzniknout tím, že je paměť alokována v jedné funkci na řádku 20, přerušena výjimkou na řádku 80 a uvolnění voláno teprve na řádku 150 -- přičemž cesta výjimky ho obejde.
    Pokud model nemá v kontextu celý tento úsek najednou, může lokálně popsat alokaci jako správnou, aniž by postřehl problém v celkové logice.
    Výzkum navíc ukazuje, že i v případech, kdy se relevantní informace do kontextového okna vejdou, modely nejspolehlivěji zpracovávají informace umístěné na začátku nebo konci vstupu -- informace uprostřed dlouhého kontextu mohou být přehlédnuty~\cite{liu-lost}.

    \item Smysluplná analýza kódu vyžaduje nejen samotný kód, ale i zadání úlohy a systémový prompt.
    Celková délka promptu tak může přesáhnout limit i u kódu, jehož samotná délka by byla přijatelná.
    Samotná délka vstupu přitom negativně ovlivňuje kvalitu výstupu modelu, a to i tehdy, když všechny relevantní informace jsou formálně v kontextovém okně přítomny~\cite{hosseini-efficient}.
\end{enumerate}

\subsubsection*{Strategie pro práci s omezeným kontextem}
\label{subsubsec:llm-token-limit-strategies}

V situacích, kdy vstup přesahuje kontextové okno modelu, se v praxi používají tři základní strategie, každá s vlastními trade-offs.

\begin{enumerate}
    \item \textbf{Chunking} -- rozdělení kódu na menší části, které jsou analyzovány samostatně.
    Výhodou je jednoduchost implementace, ale nevýhodou je ztráta globálního kontextu.
    Model při analýze každé části neví, co se děje v ostatních částech, a nemůže detekovat problémy vyplývající ze vzájemných závislostí.

    \item \textbf{Hierarchická analýza} -- nejprve jsou samostatně analyzovány menší části kódu (například jednotlivé funkce), výsledky jsou pak agregovány do celkového shrnutí.
    Tento přístup lépe zachovává globální obraz, avšak agregační krok může přenášet nebo zvyšovat nepřesnosti z dílčích analýz.
    Zároveň přidává výpočetní náklady ve formě více průchodů modelem.

    \item \textbf{Selektivní kontext} -- modelu jsou předány pouze části kódu označené jako relevantní pro hodnocenou úlohu, například konkrétní funkce, jejíž implementaci zadání specifikuje.
    Tento přístup je efektivní, avšak silně závislý na kvalitě výběru vstupu.
    Pokud výběr vynechá důležitou část, model generuje neúplnou nebo nepřesnou analýzu.
\end{enumerate}

Pro typické studentské úlohy v systému Kelvin, kde se rozsah kódu pohybuje v řádu stovek řádků, jsou všechny tři strategie ve většině případů zbytečné.
Moderní modely s kontextovým oknem 16~000 tokenů a více totiž pojmou celé odevzdání včetně zadání a systémového promptu s dostatečnou rezervou.
Strategie jsou relevantní pro rozsáhlejší projekty nebo v případě volby modelu s menším kontextovým oknem.

\endinput